\documentclass[10pt,conference]{IEEEtran}
\usepackage[T1]{fontenc}
\usepackage[utf8]{inputenc}
\usepackage{times}
\usepackage{graphicx}
\usepackage{booktabs}
\usepackage{listings}
\usepackage{enumitem}
\usepackage{xcolor}
\usepackage{url}
\usepackage{balance}
\setlist[itemize]{leftmargin=1.2em,itemsep=0.2em,topsep=0.2em}
\setlist[enumerate]{leftmargin=1.35em,itemsep=0.2em,topsep=0.2em}
\lstset{
  basicstyle=\ttfamily\scriptsize,
  columns=fullflexible,
  keepspaces=true,
  breaklines=true,
}

\title{Final Implementation Appendix --- AI-Driven Search over PhishLLM Configurations}
\author{\IEEEauthorblockN{Nazish Baliyan, CS7602 Final Submission}}

\begin{document}
\maketitle

\begin{abstract}
This appendix documents the changes between the CS7602 midterm design
document and the final implementation, the artefacts produced by an
end-to-end run, and the lessons learned during the search loop. The
implementation is in the accompanying \texttt{phishllm\_final/} repository.
\end{abstract}

\section{Changes since the midterm}
\begin{itemize}
    \item \textbf{Larger evaluation split.} The midterm prototype shipped a
    7-sample smoke-test dataset. The final implementation ships a
    deterministic synthetic generator (\texttt{src/phishllm\_search/genset.py})
    that produces 142 samples (72 phishing / 70 benign) covering six
    phishing failure modes and seven benign categories, including two
    adversarial benign classes that were not in the midterm.
    \item \textbf{Stronger candidate schema.} The final schema is
    JSON-Schema-validated at every step (load, propose, accept). The
    midterm used a Python-dict-based contract.
    \item \textbf{Real AI proposer.} The midterm used a single
    deterministic mutation list. The final loop adds an LLM proposer that
    talks to the Anthropic Claude Messages API with a structured
    meta-prompt (\texttt{prompts/meta\_search\_prompt.txt}) and a
    deterministic heuristic fallback for offline / no-API-key runs.
    \item \textbf{Failure-bucket classifier.} The midterm tracked five
    qualitative failure modes. The final classifier tracks eight (added
    \texttt{brand\_miss} and \texttt{fusion\_miss}) and is unit-tested.
    \item \textbf{Selector and stopping rules.} A precision-floor selector,
    a Pareto frontier, and explicit stopping rules
    (\texttt{max\_rounds}, \texttt{no\_recall\_gain}, \texttt{empty\_round})
    replace the midterm's manual ranking.
    \item \textbf{Reporting pipeline.} The final loop renders five
    matplotlib plots (search trace, Pareto frontier, failure buckets per
    round, top-K confusion matrix, top-K recall with bootstrap CI) and
    four Markdown/CSV tables (top-10 leaderboard, per-round failures,
    seed comparison, baseline-vs-best). The midterm produced raw JSON
    only.
    \item \textbf{Tests.} 48 unit tests across schema, metrics, failures,
    backends, evaluator, proposer, and the search loop.
    \item \textbf{HPC.} Three Slurm scripts including an array job that
    submits one task per candidate.
\end{itemize}

\section{Headline numbers}
The full search runs end-to-end in $\approx$\,2 seconds on commodity
hardware (no GPU). Twenty candidates are evaluated across three rounds;
fifteen meet the precision floor. The best candidate dominates the
official-style baseline on every reported axis:

\begin{center}
\renewcommand{\arraystretch}{1.05}
\small
\begin{tabular}{lrrrrrr}
\toprule
\textbf{Candidate} & \textbf{P} & \textbf{R} & \textbf{F1} & \textbf{rt (s)} & \textbf{\$/1K} & \textbf{Floor} \\
\midrule
\texttt{seed\_baseline}      & 1.00 & 0.74 & 0.85 & 1.70 & 2.30 & \checkmark \\
\texttt{seed\_recall\_first} & 0.78 & 1.00 & 0.88 & 1.70 & 2.30 & $\times$ \\
\texttt{seed\_no\_validation}& 0.51 & 0.74 & 0.60 & 1.30 & 1.70 & $\times$ \\
\texttt{seed\_robust}        & 1.00 & 1.00 & 1.00 & 1.70 & 2.30 & \checkmark \\
\textbf{\texttt{round2\_thr90}} (best) & \textbf{1.00} & \textbf{1.00} & \textbf{1.00} & \textbf{0.55} & \textbf{0.50} & \checkmark \\
\bottomrule
\end{tabular}
\end{center}

\section{Search trace and Pareto frontier}
The search trace shows the best-recall-so-far (under the precision floor)
growing monotonically with the largest jump in round 1, where the
heuristic proposer moves from the baseline-class candidates ($R\!=\!0.74$)
to the robust-prompt cluster ($R\!=\!1.0$). Round 2 produces no further
recall gain and the search terminates on
\texttt{no\_recall\_gain\_under\_floor}.

The Pareto frontier over (recall up, runtime down, cost down) yields three
distinct operating points: a low-cost no-interaction point
(\texttt{round2\_thr90}: \$0.50/1K, 0.55\,s), a balanced cached-validation
point (\texttt{seed\_balanced}: \$1.80/1K, 1.35\,s), and a high-recall
robust point (\texttt{seed\_robust}: \$2.30/1K, 1.70\,s). All three meet
the precision floor; the practitioner can pick the right operating point
by their cost/latency budget.

\section{Most informative experiments}
\begin{itemize}
    \item \textbf{Popularity validation is load-bearing.} Disabling it
    (\texttt{seed\_no\_validation}) raised \texttt{alias\_false\_positive}
    from 0 to 50 events --- the largest single regression observed in any
    candidate.
    \item \textbf{URL-stem fallback needs corroboration.} The
    \texttt{brand\_recall\_v1} prompt fires URL-stem fallback
    unconditionally and produces 20 \texttt{brand\_hallucination} FPs
    on indie sites with brand-like stems. The \texttt{brand\_robust\_v1}
    prompt only fires the fallback when the URL is on a suspicious TLD or
    a hosting platform, eliminating those FPs without losing any phishing
    recall.
    \item \textbf{Cached popularity validation is almost free.} The cost
    drops from \$2.30 to \$1.00 per 1K pages with no recall loss.
    \item \textbf{Prompt-injection defence has near-zero cost.}
    \texttt{prompt\_defense=true} never harmed any operating point.
\end{itemize}

\section{Reproduction}
\begin{lstlisting}[basicstyle=\ttfamily\scriptsize]
make install
make demo                       # ~10 seconds, deterministic
make search ROUNDS=4 SEED=7     # full search (heuristic proposer)
make report                     # plots + tables + case study

# To switch to the LLM proposer:
export ANTHROPIC_API_KEY=...
make search PROPOSER=llm
\end{lstlisting}

\section{Honest limitations}
The mock backend is a transparent rule emulator, not the real PhishLLM
pipeline. It is intentionally simple so the search trade-offs are
legible, and it reproduces the failure modes the paper catalogues, but it
does not reproduce every nuance of the upstream OCR / captioning / Google
Programmable Search integration. The included \texttt{OfficialRepoBackend}
documents the exact contract for plugging in the upstream repository; the
search loop, evaluator, reporting, and Slurm scripts are unchanged.

The headline numbers are reported on a synthetic 142-site reduced split.
A full-scale evaluation on the 12K-sample public benchmark is feasible
under the same contract but was out of scope for the 5-day final-submission
window.

\end{document}
